Negli ultimi anni l'intelligenza artificiale ha assunto un ruolo sempre più centrale nella nostra quotidianità, grazie soprattutto alla rapida diffusione dei Large Language Model (LLM). L'utilizzo di questi modelli sta progressivamente evolvendo: da semplici strumenti di generazione testuale a componenti centrali di sistemi software in grado di pianificare, decidere e agire autonomamente. Nasce così il concetto di agente. Questa evoluzione ha aperto la strada a scenari applicativi in cui più agenti possono comunicare e collaborare tra loro al fine di portare a termine compiti complessi.

In un contesto in cui gli agenti comunicano tra loro nasce la necessità di introdurre protocolli di comunicazione standardizzati che ne facilitino l'interazione. Gli agenti devono essere in grado di scoprirsi reciprocamente e scambiarsi informazioni, in modo da potersi coordinare in attività condivise indipendentemente dalla natura implementativa di ciascuno di essi. Il protocollo Agent2Agent (A2A) nasce proprio per rispondere a questa esigenza, definendo un insieme di regole e strutture di dati -- Message, Task, Artifact, Agent Card -- che sono alla base di ogni interazione tra agenti.

La tesi nasce nell'ambito di un tirocinio curriculare svolto presso l'azienda HCL Software, il cui obiettivo era progettare e implementare la gestione del ciclo di vita dei Task all'interno del protocollo A2A, utilizzando framework come LangGraph per la costruzione degli agenti. L'attenzione si è concentrata principalmente sulla classe Python \texttt{AgentExecutor}, e più precisamente sui suoi due metodi principali, \texttt{execute} e \texttt{cancel}, responsabili rispettivamente dell'esecuzione e dell'interruzione di un Task.

L'obiettivo di questa tesi è duplice: da un lato fornire una comprensione approfondita del protocollo A2A e del ruolo che la classe \texttt{AgentExecutor} svolge al suo interno; dall'altro illustrare il percorso di progettazione e implementazione che ha portato alla realizzazione di un sistema in grado di gestire correttamente i Task, anche in scenari complessi come la cancellazione di un Task durante una comunicazione in streaming.

Il contributo pratico di questo lavoro consiste nell'implementazione completa dei metodi \texttt{execute} e \texttt{cancel} in una classe comune, applicabile ad agenti diversi basati su LangGraph, nella realizzazione di un server A2A funzionante e nella validazione del sistema attraverso una serie di scenari di test mirati.

La tesi è organizzata come segue. Il capitolo 1 introduce la distinzione tra un LLM e un agente. Il capitolo 2 descrive il protocollo A2A, motivandone la nascita e approfondendo l'architettura, il ciclo di vita dei Task e il ruolo della classe \texttt{AgentExecutor}. Il capitolo 3 espone i requisiti del sistema da realizzare. Il capitolo 4 presenta la progettazione dei metodi \texttt{execute} e \texttt{cancel}. Il capitolo 5 descrive l'implementazione realizzata. Il capitolo 6 illustra la validazione del sistema attraverso i test svolti. Infine, le conclusioni riassumono il lavoro svolto e delinea possibili sviluppi futuri.